\section{Yêu cầu đề tài}

\subsection{Bối cảnh}

Cơ chế attention là thành phần cốt lõi của Transformer. Trong mỗi lớp attention, điểm số giữa các vị trí được chuẩn hoá theo hàng bằng hàm softmax trước khi nhân với giá trị. Với độ dài chuỗi $L$, mỗi hàng cần $L$ phép tính $e^{x}$, $L-1$ phép cộng và một phép chia, và điều này lặp lại theo số hàng và số đầu attention. Khi mô hình được lượng tử hoá sang INT8, phần tính softmax trở thành điểm nghẽn nếu vẫn dùng số thực dấu phẩy động.

\subsection{Mục tiêu}

Đề tài xây dựng một bộ tăng tốc softmax theo hàng, trong đó toàn bộ phép tính dùng số nguyên, nhận đầu vào INT8 có dấu và trả về đầu ra UINT8 với scale $1/256$. Bộ tăng tốc được hiện thực dưới dạng IP Avalon-MM slave trên kit DE10-Standard (Cyclone V 5CSXFC6D6F31C6), dùng Quartus Prime Lite 18.1.

Các mục tiêu cụ thể:
\begin{itemize}
  \item Số học bit-exact so với một golden model viết bằng Python.
  \item Tần số hoạt động từ 100 MHz trở lên sau Fitter.
  \item Tổng số ALUT và DSP vượt 1000 theo yêu cầu môn học.
  \item Chạy được trên board qua JTAG và System Console, không cần Linux.
\end{itemize}

\subsection{Yêu cầu chi tiết}
\label{sec:req}

Các yêu cầu được chốt ở phiên bản 1.0 của tài liệu yêu cầu và là hợp đồng giao diện giữa các nhóm; Bảng \ref{tab:req} liệt kê các yêu cầu chính cùng mã yêu cầu. Cấu hình mặc định là $P = 8$ lane, $N_{\text{SEG}} = 8$ đoạn PWL, $L_{\max} = 512$, $R_{\max} = 16$, và các tham số này đọc được qua thanh ghi \texttt{HW\_PARAM}.

\begin{table}[htbp]
\centering
\small
\begin{tabular}{p{2.0cm}p{5.2cm}p{6.6cm}}
\toprule
Mã & Yêu cầu & Chỉ tiêu \\
\midrule
FUN-01 đến 03 & Tính softmax theo hàng & $1 \le L \le 512$, $1 \le R \le 16$; $L$ không cần là bội của $P$; phần tử $i \ge L$ không ảnh hưởng kết quả \\
FUN-04, 05 & Cấu hình sai và độc lập lần chạy & $L$ hoặc $R$ ngoài khoảng: không chạy, DONE và ERR lên; kết quả chỉ phụ thuộc $X$, $L$, $K$ \\
ARI-01 đến 08 & Số học bit-exact & Vào INT8, ra UINT8 với scale $1/256$; $K$ dạng UQ0.16 do phần mềm tính; PWL $N_{\text{SEG}}$ đoạn cho $2^F$; hai lượt quét với IntMax; một phép chia mỗi hàng; bão hoà về 255 \\
IF-01 đến 03 & Giao diện bus & Avalon-MM slave 32 bit, địa chỉ theo word, không burst, không \texttt{waitrequest}, độ trễ đọc 2 chu kỳ, \texttt{byteenable} 4 bit \\
IF-05, 10 & Báo hoàn thành & Thăm dò bit DONE, không ngắt; conduit \texttt{status\_busy}, \texttt{status\_done}, \texttt{status\_err} \\
PRF-01 & Tần số & $\ge 100$ MHz trên Cyclone V 5CSXFC6D6F31C6 \\
PRF-03 & Độ trễ mỗi hàng & $\le 2\lceil L/P \rceil + 64$ chu kỳ, tức $\le 192$ chu kỳ ở $L = 512$, $P = 8$ \\
RES-01 & Tài nguyên & Tổng ALUT và DSP lớn hơn 1000 sau Fitter \\
IMP-01 đến 04 & Hiện thực & SystemVerilog tập con Quartus Prime Lite 18.1; bảng PWL sinh từ golden model, không sửa tay \\
VER-01 đến 04 & Kiểm thử & So bit-exact với golden model; danh sách kịch bản tối thiểu; khẳng định giao thức; \texttt{CYCLE\_CNT} theo công thức độ trễ \\
\bottomrule
\end{tabular}
\caption{Các yêu cầu chính của đề tài, theo tài liệu yêu cầu phiên bản 1.0.}
\label{tab:req}
\end{table}

\subsection{Phạm vi}

Đây là báo cáo tiến độ ở mốc 50\% của đề tài: lõi đã được thiết kế, kiểm thử bằng mô phỏng và tổng hợp riêng. Tích hợp hệ thống, chạy trên board và đánh giá trên mô hình mục tiêu thuộc giai đoạn còn lại; các ô ghi [Chưa đo] sẽ được bổ sung khi đó. Báo cáo tập trung vào thiết kế lõi \texttt{smx\_core}, gồm thanh ghi điều khiển, bộ đệm X/Y, bộ điều khiển và datapath. Lớp bao Avalon, hệ thống Platform Designer và script demo do nhóm hiện thực đảm nhiệm; báo cáo mô tả các phần này ở mức giao diện và cấu trúc hệ thống (mục \ref{sec:iface}). Phần kiểm thử lõi do nhóm kiểm thử thực hiện qua cổng \texttt{acc\_*}; kế hoạch và kết quả được trình bày ở Phần 3. Hình \ref{fig:smx-system} cho thấy vị trí của lõi trong hệ thống demo và ranh giới giữa các phần việc.

\begin{figure}[htbp]
  \centering
  \includegraphics[width=0.84\textwidth]{smx_system.pdf}
  \caption{Ranh giới hệ thống: PC, cầu JTAG to Avalon Master, lớp bao \texttt{smx\_avalon} và lõi \texttt{smx\_core}. Phần trong khung nét đứt nằm trên FPGA.}
  \label{fig:smx-system}
\end{figure}

Đề tài không bao gồm phép nhân ma trận, convolution, LayerNorm, GELU và masking. Hàm sigmoid và GELU chỉ được xem là mở rộng tuỳ chọn sau khi lõi hoàn thành.

\subsection{Đóng góp}

Các đóng góp chính của đề tài gồm:
\begin{enumerate}
  \item Thuật toán softmax thuần số nguyên, không dùng hàm log trong datapath, nhờ đưa hệ số $\log_2 e$ vào $K$ do phần mềm tính sẵn.
  \item Kết hợp IntMax với chuẩn hoá trực tuyến để chỉ cần hai lượt quét và hiệu chỉnh tổng bằng phép dịch.
  \item Tính lại $E$ ở lượt thứ hai từ bộ đệm đầu vào, thay vì lưu kết quả trung gian.
  \item Golden model bit-exact và đặc tả giao diện, làm cơ sở để các nhóm khác tích hợp và kiểm thử.
  \item Lõi RTL đã kiểm thử bit-exact với golden model và đạt mục tiêu tần số, tài nguyên sau Fitter.
\end{enumerate}

\subsection{Tóm tắt kết quả}

Bảng \ref{tab:summary} tóm tắt các kết quả chính của đề tài ở cấu hình mặc định; chi tiết và điều kiện đo nằm ở các mục được dẫn trong bảng.

\begin{table}[htbp]
\centering
\small
\begin{tabular}{lll}
\toprule
Hạng mục & Kết quả & Mục \\
\midrule
Cấu hình mặc định & $P = 8$ lane, $N_{\text{SEG}} = 8$, $L_{\max} = 512$, $R_{\max} = 16$ & \ref{sec:req}, \ref{sec:flow} \\
Sai số so với softmax float64 & khoảng 0,002, bằng nửa LSB của đầu ra & \ref{sec:err} \\
Tài nguyên sau Fitter & 1\,681 ALUT, 24 DSP, 21 khối M10K & \ref{sec:synth} \\
Tần số tối đa & 122 đến 134 MHz qua ba seed & \ref{sec:synth} \\
Độ trễ một hàng, $L = 512$ & 182 chu kỳ & \ref{sec:sim-res}, \ref{sec:synth} \\
Kiểm thử lõi & 7 testcase, 458\,837 phép so, khớp bit-exact & \ref{sec:sim-res} \\
Chạy trên board & [Chưa đo] & \ref{sec:plan} \\
\bottomrule
\end{tabular}
\caption{Tóm tắt kết quả chính ở cấu hình mặc định.}
\label{tab:summary}
\end{table}

\subsection{Cấu trúc báo cáo}

Báo cáo gồm ba phần theo yêu cầu và một phần về tình hình thực hiện. Phần 1 là phần này. Phần 2 mô tả thiết kế: mục \ref{sec:principle} trình bày nguyên lý hoạt động và mục \ref{sec:rtl} trình bày sơ đồ RTL của lõi và của từng khối. Phần 3 trình bày kiểm thử bằng mô phỏng và kết quả tổng hợp. Phần 4 nêu tình hình thực hiện và kế hoạch cho giai đoạn còn lại.
